\documentclass[a4paper]{article}

\usepackage[a4paper, top=8mm, bottom=8mm, left=8mm, right=8mm]{geometry}

\usepackage{polyglossia}
\setdefaultlanguage[babelshorthands=true]{russian}
\setotherlanguage{english}

\usepackage{fontspec}
\setmainfont{FreeSerif}
\newfontfamily{\russianfonttt}[Scale=0.7]{DejaVuSansMono}

\usepackage[tiny, compact]{titlesec}

\usepackage{titling}
\setlength{\droptitle}{-1cm}
\pretitle{\begin{center}\begin{bfseries}\Large}
\posttitle{\par\end{bfseries}\end{center}}
\preauthor{\begin{center}\normalsize}
\postauthor{\par\end{center}\vspace{-1.8cm}}

\usepackage{hyperref}
\usepackage{bookmark}
\usepackage{csquotes}

\title{UCFS: Доработка репозитория}

\author{Исаев Даниил Сергеевич}

\date{}

\begin{document}

\maketitle

\begin{flushright}
    Группа:~\emph{25.Б72-мм}

    Кафедра:~\emph{Кафедра системного программирования СПбГУ}

    Научный руководитель:~\emph{Григорьев Семён Вячеславович}
    
    Номер семестра практики:~\emph{2}
\end{flushright} 

\section{Введение}

Работа выполняется в рамках проекта UCFS\footnote{Ссылка на UCFS: \url{https://github.com/FormalLanguageConstrainedPathQuerying/UCFS}}. UCFS~--- это универсальный инструмент, предназначенный для решения задач на стыке теории контекстно-свободных языков и анализа ориентированных графов с метками на рёбрах. В~основе его работы лежит алгоритм GLL~(Generalized LL parsing), обобщающий классический LL-анализ на случай неоднозначных грамматик, а также его расширение для работы с графами. Этот инструмент применим для широкого круга задач, включая контекстно-свободные запросы к путям в графах~(Context-Free Path Querying, CFPQ) и задачу контекстно-свободной достижимости~(Context-Free Language Reachability, CFL-R).

На данный момент проект обладает широкой функциональностью, однако существующие примеры использования опираются на сложные прикладные грамматики и большие графы из реальных задач, что затрудняет понимание базовых принципов работы новыми пользователями. Данная практика направлена на улучшение документации и создание пошаговой инструкции, включающей простые примеры для запуска, что позволит быстрее освоить UCFS и начать применять его в учебных и исследовательских задачах.

\section{Постановка задачи}

\textbf{Целью работы является} создание пошаговой инструкции по аналогии с популярными библиотеками, а также простых примеров использования инструмента UCFS, включая реализацию кода и интеграцию инструкции по запуску в основную документацию.

\textbf{Для достижения цели требуется решить следующие задачи.}
\begin{enumerate}
    \item Изучить материалы, касающиеся контекстно-свободных грамматик и языков.
    \item Выполнить обзор существующих примеров UCFS и подходов к написанию документации.
    \item Спроектировать формат интеграции учебных примеров и структуру документации для быстрого старта.
    \item Реализовать набор простых учебных примеров на языке Kotlin с использованием встроенного DSL, подготовить описания графов в формате~\texttt{.dot} и их визуализации.
    \item Выполнить тестирование реализованных примеров на корректность работы.
    \item Выполнить реструктуризацию документации проекта: оставить README-файлы техническими, а учебную и справочную информацию вынести в специализированные разделы документации.
    \item Оформить полученные результаты в виде пуллреквеста.
\end{enumerate}

\section{Обзор}

В рамках проекта UCFS уже существовал набор примеров использования. Однако анализ текущей документации показал, что существующие примеры оперируют сложными грамматиками и большими графами, взятыми из реальных прикладных задач. Это затрудняет понимание базовых принципов работы инструмента пользователями, которые только начинают знакомство с CFPQ и GLL-анализом.

В других проектах, таких как NetworkX\footnote{Ссылка на документацию NetworkX: \url{https://networkx.org/documentation/stable/}}, стандартной практикой является наличие разделов \enquote{Tutorial} с инструкцией по работе с инструментом и \enquote{Getting Started} с минимальными воспроизводимыми примерами, демонстрирующими базовый API. В данном проекте вышеперечисленные разделы отсутствуют.

Таким образом, для эффективного освоения инструмента новыми пользователями необходимо создать упрощённые учебные сценарии с пошаговыми инструкциями, аналогичные тем, что используются в популярных проектах.

\section{Описание предлагаемого решения}

Предлагаемое решение заключается в создании документации, которая пошагово направляет нового пользователя при начале работы с UCFS. Она включает базовую информацию о контекстно-свободных языках~(Context-Free Languages, CFL), инструкции по заданию графов и грамматик, руководство по интерпретации выходных данных, а также набор минимальных примеров для последовательного знакомства с возможностями инструмента. В отличие от существующих сложных сценариев, разработанные примеры намеренно упрощены: графы имеют небольшой размер, а их структура интуитивно понятна. В качестве ключевого примера выбрана контекстно-свободная грамматика $a^nb^n$. Являясь каноническим примером, она легко воспринимается начинающими и широко используется в учебных пособиях по формальным языкам для иллюстрации базовых концепций. Это позволяет сосредоточиться на изучении возможностей инструмента, не отвлекаясь на специфику сложных предметных областей. Для описания грамматики используется встроенный DSL проекта UCFS, а графы представлены в формате~\texttt{.dot}.

Принятые решения обоснованы следующим образом: простота примеров снижает порог входа и подготавливает новых пользователей к работе со сложными графами, имеющими практическое применение. Интеграция инструкции по запуску непосредственно в основной файл документации (вместе с визуализациями графов в формате SVG) устраняет необходимость поиска отдельной информации.

Реализация решения осуществлялась с использованием языка Kotlin~(код примеров), DSL для задания грамматики, формата~\texttt{.dot} для описания графов, консольной утилиты dot~(Graphviz) для конвертации графов в SVG. Для оформления материалов использовалась система MkDocs с разметкой Markdown, что позволило создать структурированный и удобный для навигации справочный ресурс.

\section{Тестирование}

Проверка разработанных примеров проводилась путем ручной верификации структуры графов с последующим построением SPPF~(Shared Packed Parse Forest) с помощью инструмента UCFS.
Критериями успешного тестирования являлись: успешная компиляция и запуск всех примеров, а также полное совпадение построенных лесов разбора с теоретически ожидаемыми результатами. В ходе проверки все примеры успешно прошли компиляцию, а анализ полученных SPPF подтвердил их соответствие теоретическим ожиданиям для выбранных грамматик и графов, что свидетельствует о корректности реализации.

\section{Заключение}

В ходе выполнения работы получены следующие результаты.

\begin{itemize}
    \item Реализован набор минимальных учебных примеров на языке Kotlin с использованием встроенного DSL, подготовлены исходные описания графов в формате~\texttt{.dot} и их итоговые визуализации.
    \item Разработана пошаговая инструкция для быстрого старта и выполнена реструктуризация документации проекта с чётким разделением технического README и учебных материалов.
    \item Получены результаты тестирования, подтверждающие корректность работы примеров.
    \item Разработчиками проекта принят пуллреквест, содержащий весь код, визуализации и обновлённую документацию.
\end{itemize}

Репозиторий (форк): \url{https://github.com/danisaev/UCFS}.

Ссылка на пуллреквест: \url{https://github.com/vkutuev/UCFS/pull/3}.

\end{document}
